
\hspace{0.5cm}P1) \textcolor{red}{[Control 3 - 2024-20]} Un cilindro de masa m y radio R (momento de inercia Icm) es forzado con una tensi\'on T aplicada a trav\'es de un cuerda enrollada a un carrete de radio r adosado al cilindro. Asumiendo que el cilindro gira sin deslizar sobre una superficie horizontal rugosa:\\
\begin{enumerate}[a)]
\hspace{2cm}\item Escriba la condici\'on de rodadura entre la velocidad de traslaci\'on y la velocidad de rotaci\'on. Derive y deduzca la relaci\'on para las aceleraciones.
\hspace{2cm}\item Realice DCL de cilindro
\hspace{2cm}\item Escriba las ecuaciones de movimiento
\hspace{2cm}\item Calcule aceleraci\'on a del centro de cilindro
\hspace{2cm}\item Asumiendo que el cilindro parte del reposo, calcule la velocidad del cilindro v cuando \'este haya dado una vuelta, en funci\'on de los datos del problema
\end{enumerate}




